\chapter{Pipeline de IA}

O pipeline de IA foi separado em componentes independentes para permitir que cada etapa possa ser evoluída sem afetar as demais. No estado atual, os componentes existem como placeholders arquiteturais, preparados para receber a implementação real quando o modelo treinado estiver disponível. A intenção é manter a sequência de processamento clara, rastreável e substituível.

\section{Fluxo de processamento}

O fluxo lógico do pipeline é o seguinte:
\begin{center}
Preprocessor $\rightarrow$ ModelLoader $\rightarrow$ ModelRunner $\rightarrow$ Prediction $\rightarrow$ PredictionPostProcessor $\rightarrow$ PredictionResult
\end{center}

O \textit{ImagePreprocessor} é responsável por transformar a imagem em uma representação numérica compatível com o modelo. O \textit{ModelLoader} abstrai o carregamento do modelo e sua permanência em memória. O \textit{ModelRunner} executará a predição futura usando o modelo carregado. A estrutura \textit{Prediction} representa a saída bruta do modelo, enquanto o \textit{PredictionPostProcessor} converte essa saída em \textit{PredictionResult}, já no formato consumível pela aplicação.

\section{Responsabilidades de cada etapa}

O \textit{ImagePreprocessor} deve cuidar de redimensionamento, normalização e transformação do fluxo de imagem para tensor ou vetor numérico. O \textit{ModelLoader} será responsável por carregar o artefato treinado apenas uma vez e manter uma instância reutilizável em memória. O \textit{ModelRunner} utilizará o modelo carregado para realizar inferência. O \textit{PredictionPostProcessor} aplica reestruturação semântica ao resultado bruto, produzindo um objeto final mais adequado ao domínio da aplicação.

\section{Desacoplamento}

A separação por interfaces permite que o pipeline seja evoluído de forma incremental. Em vez de depender de detalhes de TorchSharp, PyTorch ou ONNX diretamente em toda a aplicação, apenas o componente responsável pela execução do modelo será alterado quando a implementação real chegar. Essa abordagem reduz o impacto de alterações na tecnologia de inferência.

\begin{figure}[H]
    \centering
    \fbox{\parbox{0.85\textwidth}{\centering Espaço reservado para o diagrama do pipeline de IA.}}
    \caption{Pipeline de IA do microserviço}
    \label{fig:pipeline-ia}
\end{figure}
